

\chapter{Erwartete Ergebnisdarstellung}
Dieses Kapitel beschreibt, wie die Ergebnisse später dargestellt werden.

Zunächst werden die Prognosefehler auf den unveränderten Testdaten gegenübergestellt. Danach sollen Diagramme zeigen, wie sich die Ergebnisse bei zunehmenden Datenfehlern verändern. Die Streuung zwischen den Versuchsläufen und geeignete Konfidenzintervalle helfen dabei, die Unterschiede einzuordnen. Ergebnisse der beiden Teildatensätze werden getrennt dargestellt.

Die Messungen zum Cloud-Betrieb werden in Tabellen und Diagrammen zusammengefasst. Ein Vergleich von Prognosefehlern und Betriebskosten soll zeigen, ob eine genauere Prognose mit zusätzlichem Aufwand verbunden ist. Daraus soll eine Einschätzung der Stärken und Schwächen der verschiedenen Modelle unter Berücksichtigung der untersuchten Bedingungen entstehen.
